
\documentclass{article}
\usepackage{colortbl}
\usepackage{makecell}
\usepackage{multirow}
\usepackage{supertabular}

\begin{document}

\newcounter{utterance}

\centering \large Interaction Transcript for game `cladder', experiment `full\_v1.5\_default', episode 5573 with Qwen3.5{-}4B.
\vspace{24pt}

{ \footnotesize  \setcounter{utterance}{1}
\setlength{\tabcolsep}{0pt}
\begin{supertabular}{c@{$\;$}|p{.15\linewidth}@{}p{.15\linewidth}p{.15\linewidth}p{.15\linewidth}p{.15\linewidth}p{.15\linewidth}}
    \# & \multicolumn{2}{c}{Player} && \multicolumn{2}{c}{Game Master} \\
    \hline

    \theutterance \stepcounter{utterance}  
    & & & \multicolumn{4}{p{0.6\linewidth}}{
        \cellcolor[rgb]{0.9,0.9,0.9}{
            \makecell[{{p{\linewidth}}}]{
                \texttt{\tiny{[P1$\langle$GM]}}
                \texttt{You are an expert in causal inference. The following question is not a typical commonsense query, but rather a meticulously designed question created by a professor specializing in causal inference, intended to assess the students' mastery of the course content.} \\
\texttt{Imagine a self{-}contained, hypothetical world with only the following conditions, and without any unmentioned factors or causal relationships: Gender has a direct effect on treatment and peanut allergy. Treatment has a direct effect on peanut allergy. For patients who are not male and not receiving treatment, the probability of being allergic to peanuts is 98\%. For patients who are not male and receiving treatment, the probability of being allergic to peanuts is 87\%. For patients who are male and not receiving treatment, the probability of being allergic to peanuts is 34\%. For patients who are male and receiving treatment, the probability of being allergic to peanuts is 14\%. The overall probability of male gender is 43\%. For patients receiving treatment, would it be less likely to see being allergic to peanuts if the patient had received no treatment?} \\
\texttt{Start your answer with 'Yes' or 'No', followed by additional reasoning or evidence to support your explanation.} \\
            }
        }
    }
     \\ \\

    \theutterance \stepcounter{utterance}  
    & \multicolumn{4}{p{0.6\linewidth}}{
        \cellcolor[rgb]{0.561,0.737,0.561}{
            \makecell[{{p{\linewidth}}}]{
                \texttt{\tiny{[P1$\rangle$GM]}}
                \texttt{problem} \\
\texttt{Game Master: You are an expert in causal inference. The following question is not a typical commonsense query, but rather a meticulously designed question created by a professor specializing in causal inference, intended to assess the students' mastery of the course content.} \\
\texttt{Imagine a self{-}contained, hypothetical world with only the following conditions, and without any unmentioned factors or causal relationships: Gender has a direct effect on treatment and peanut allergy. Treatment has a direct effect on peanut allergy. For patients who are not male and not receiving treatment, the probability of being allergic to peanuts is 98\%. For patients who are not male and receiving treatment, the probability of being allergic to peanuts is 87\%. For patients who are male and not receiving treatment, the probability of being allergic to peanuts is 34\%. For patients who are male and receiving treatment, the probability of being allergic to peanuts is 14\%. The overall probability of male gender is 43\%. For patients receiving treatment, would it be less likely to see being allergic to peanuts if the patient had received no treatment?} \\
\texttt{Start your answer with 'Yes' or 'No', followed by additional reasoning or evidence to support your explanation.} \\
\\ 
\texttt{problem} \\
\texttt{RULES: } \\
\texttt{1. Start your response with the word 'problem'.} \\
\texttt{2. Provide a complete, mathematically rigorous derivation based on the provided data.} \\
\texttt{3. Explicitly calculate the counterfactual probability \$P(\textbackslash{}text\{Allergy\} | \textbackslash{}text\{Treatment\}=1, \textbackslash{}text\{No Treatment\})\$ by leveraging the causal structure (\$G \textbackslash{}rightarrow T\$) and the Law of Total Probability.} \\
\texttt{4. Compare the counterfactual probability with the observed probability \$P(\textbackslash{}text\{Allergy\} | \textbackslash{}text\{Treatment\}=1)\$ to determine if the risk decreases upon removing treatment.} \\
\texttt{5. Conclude with a definitive 'Yes' or 'No' supported by the calculated values.} \\
\\ 
\texttt{STRATEGY: } \\
\texttt{1. Define the variables: \$G\$ (Gender, 1=Male, 0=Female), \$T\$ (Treatment, 1=Treated, 0=Untreated), \$A\$ (Allergy).} \\
\texttt{2. List the known conditional probabilities: \$P(A|G=0, T=0)=0.98\$, \$P(A|G=0, T=1)=0.87\$, \$P(A|G=1, T=0)=0.34\$, \$P(A|G=1, T=1)=0.14\$.} \\
\texttt{3. Use the marginal probability \$P(G=1)=0.43\$ to establish \$P(G=0)=0.57\$.} \\
\texttt{4. Recognize that the specific proportion of males/females *within* the treated group is unknown, but the question asks about the counterfactual for the *treated* group. However, without knowing the treatment propensity scores (\$P(T=} \\
            }
        }
    }
    & & \\ \\

    \theutterance \stepcounter{utterance}  
    & & & \multicolumn{2}{p{0.3\linewidth}}{
        \cellcolor[rgb]{0.9,0.9,0.9}{
            \makecell[{{p{\linewidth}}}]{
                \texttt{\tiny{[GM$|$GM]}}
                \texttt{ParseError: The response doesn't start with one of the required word {[}'yes', 'no'{]}, but problem} \\
            }
        }
    }
    & & \\ \\

    \theutterance \stepcounter{utterance}  
    & & & \multicolumn{2}{p{0.3\linewidth}}{
        \cellcolor[rgb]{0.9,0.9,0.9}{
            \makecell[{{p{\linewidth}}}]{
                \texttt{\tiny{[GM$|$GM]}}
                \texttt{game\_result = ABORT} \\
            }
        }
    }
    & & \\ \\

\end{supertabular}
}

\end{document}
